%%%%%%%%%%%%%%%%%%%%%%%%%%%%%%%%%%%%%%%%%%%%%%%%%%%%%%%%%%%%%%%%%%%%%%%%%%%%%%%%
%% 
\cleardoubleoddpage%  Make sure to start each chapter on a new odd page
\chapter{Methodology}
\label{sec:methodology}


\section{Analysis Environment}
\label{sec:methodology:installation}
The analyses are conducted in a dedicated 64-bit virtual machine running Debian 13.6 with the Xfce desktop environment. The virtual machine is hosted on a 64-bit system using 
Oracle VirtualBox version 7.2.8 and is allocated 4096 MB of memory and four virtual processor cores. Using the same virtualized environment for all evaluated tools 
provides consistent operating-system and resource conditions, thereby reducing environmental differences that could influence the comparison.

Network access is initially available to download and install the analysis tools from their official sources. Shared folders or other VirtualBox integration features may also be 
enabled temporarily to transfer the required samples to the guest system. Once the tools and samples have been transferred, networking and all integration features that permit 
communication between the host and guest systems are disabled. These features include shared folders, the shared clipboard, and drag and drop. A snapshot of the prepared and isolated 
environment is then created, providing a consistent initial state and allowing the virtual machine to be restored if necessary.

The samples are examined exclusively through static analysis and are never intentionally executed. Nevertheless, potentially harmful files remain confined to the virtual machine 
while the analyses are performed. The virtualized environment and the disabled communication features reduce the risks associated with handling such files, although the virtual 
machine is treated as a risk-reduction measure rather than an absolute security boundary.


\section{Test Dataset}
\label{sec:methodology:dataset}


\subsection{Synthetic Samples}
\label{sec:methodology:dataset:synthetic}
Three synthetic samples are used in the evaluation. They consist of equivalent implementations of an IBAN validator written in C, Java, and C\#. 
The samples perform the same principal operations: they normalize the supplied IBAN, verify its general format and country-specific length, and apply 
the modulo-97 checksum defined for IBAN validation. The programs determine only whether an IBAN is structurally and mathematically valid.

Using implementations of the same program reduces functional differences between the samples while allowing tools for native code, Java bytecode, 
and .NET assemblies to be evaluated using comparable program logic. The implementations contain multiple functions or methods, conditional branches, 
loops, string processing, constants, and structured data. These elements provide a known basis for assessing how accurately the analysis tools reconstruct 
the original program structure and logic.

The C implementation is compiled without optimization and without debug information. The resulting executable is additionally stripped, removing symbol 
information that could reveal original function or variable names. Consequently, the native-code analysis tools must infer the program structure primarily 
from the machine code. The Java and C\# implementations are compiled to their respective intermediate representations and analyzed using Bytecode Viewer and ILSpy.

The complete source code of all three synthetic samples is provided in Appendix~\ref{app:an-appendix}.


\subsection{Real Malware}
\label{sec:methodology:dataset:real}
In addition to the synthetic samples, a real-world malware sample from the WannaCry ransomware family is used in the evaluation. The selected sample is a 32-bit 
Windows Portable Executable identified by the following SHA-256 hash:

\begin{center}
\texttt{24d004a104d4d54034dbcffc2a4b19a11f39008a575aa614ea04703480b1022c}
\end{center}

A native-code executable was selected because it can be examined by Ghidra, Binary Ninja, and Cutter. Consequently, the same sample can be analysed with three of the 
evaluated tools, providing a broader basis for comparison than a Java bytecode or .NET sample. The specimen is not packed and contains numerous elements relevant to 
static analysis, including imported functions, embedded strings, resources, function boundaries, and control-flow structures. Furthermore, WannaCry is extensively documented, 
allowing observations made during the analysis to be compared with previously established information about the malware.

The objective of the evaluation is not to reverse engineer the complete malware or provide a comprehensive account of its behaviour. Instead, the analysis concentrates on 
selected and consistently observable elements, such as executable metadata, imports, strings, cross-references, identified functions, control-flow graphs, and the decompilation 
of selected functions. This limited scope permits the output and usability of the three native-code analysis tools to be compared under equivalent conditions.

The sample is examined exclusively through static analysis and is never intentionally executed. Its cryptographic hash is stated to identify the selected specimen unambiguously 
without relying solely on its filename.


\section{Evaluation Procedure}
\label{sec:methodology:procedure}
The evaluation follows a standardized procedure to ensure that the applicable tools are examined under comparable conditions. The native C sample is analysed using Ghidra, 
Binary Ninja, and Cutter, while the Java and C\# samples are examined using Bytecode Viewer and ILSpy, respectively. The WannaCry sample is analysed with Ghidra, Binary Ninja, and Cutter. 
Each tool receives the same applicable input file, and none of the samples are executed at any stage of the evaluation.

For each tool, the relevant sample is imported and analysed using the default settings provided by the tool. If additional information, such as the target architecture or 
file format, must be specified manually, the selected configuration is documented. The automatically generated results are examined before function names, variable names, 
data types, or comments are modified manually. This ensures that the evaluation reflects the output initially produced by each tool rather than improvements introduced by the analyst.

The examination concentrates on a consistent set of program elements. These include file-format recognition, detected functions or methods, imported functions and libraries, 
embedded strings and constants, cross-references, control-flow representations, and decompiled output. For the synthetic samples, the reconstructed output is compared with the 
corresponding original source code. This provides a known reference for assessing the preservation of program structure, control flow, and logic. The findings obtained from 
the WannaCry sample are compared across the three native-code analysis tools and, where appropriate, with publicly documented information about the malware.

Observations are recorded according to the evaluation criteria defined in Chapter~\ref{sec:evaluation}. Relevant screenshots are captured to illustrate the user interface, analysis workflow, 
and representative differences in the generated output. The procedure is repeated for every applicable combination of tool and sample, allowing the results to be compared systematically.

%% 
%%%%%%%%%%%%%%%%%%%%%%%%%%%%%%%%%%%%%%%%%%%%%%%%%%%%%%%%%%%%%%%%%%%%%%%%%%%%%%%%
